% Gap-Serie Teil 4; korrigiertes Arbeitspapier v0.3, 2026-10-02.
% Historischer quellengebundener Formaudit; allein R19-Inventuränderungen.
% v0.2-Originale und ihre Claim-/Quellenbindungen bleiben geschützt.
\documentclass[11pt,a4paper]{article}
\usepackage{cmap}
\usepackage[utf8]{inputenc}
\usepackage[T1]{fontenc}
\usepackage[ngerman]{babel}
\usepackage{lmodern}
\usepackage{microtype}
\usepackage{amsmath,amssymb,amsthm}
\usepackage{booktabs}
\usepackage{tabularx}
\usepackage{array}
\usepackage[colorlinks=true,linkcolor=blue,citecolor=blue,urlcolor=blue]{hyperref}
\renewcommand*{\HyperDestNameFilter}[1]{de.#1}
\usepackage{xurl}
\input{glyphtounicode}
\pdfgentounicode=1

\hypersetup{
  pdftitle={Das Bridge-Zertifikat, angewandt auf die Joshi-ATS-Serie: Ein Ledger-Audit},
  pdfauthor={Lukas Geiger},
  pdfsubject={Vorab festgelegter Vier-Zeilen-Ledger-Audit der Arithmetischen Teichmüller-Raum-Serie von Kirti Joshi},
  pdfkeywords={abc-Vermutung, Inter-universelle Teichmüller-Theorie, Arithmetische Teichmüller-Räume, Korollar~3.12, Bridge-Zertifikat, Ledger-Audit},
  pdflang={de},
  bookmarksnumbered=true,
  pdfpagemode=UseOutlines
}

\clubpenalty=10000
\widowpenalty=10000
\displaywidowpenalty=10000
\frenchspacing

\newtheorem{remark}{Bemerkung}[section]

\title{Das Bridge-Zertifikat, angewandt auf die Joshi-ATS-Serie:\\Ein Ledger-Audit}
\author{Lukas Geiger\\\small Unabhängiger Forscher, Bernau, Deutschland\\\small ORCID: 0009-0005-7296-1534}
\date{2. Oktober 2026 --- Version 0.3 (Arbeitspapier)}

\begin{document}
\maketitle

\begin{abstract}
Wir überarbeiten den vorab festgelegten Formaudit von Kirti Joshis Serie
zum Arithmetischen Teichmüller-Raum (ATS), unter Beibehaltung der
Quellenversionen und der historischen Auswahl der Zeilen R2, R4, R6, R8.
Dies ist ein Formaudit ausgewählter Zeilen, keine vollständige
Zertifizierung und kein Audit der Beweiskorrektheit. Die Quelle macht das
$j^2$-Skalierungsgesetz, die zugewiesene kompensierende Stelle und den
asymmetrischen Frobenius-Shift sichtbar. Diese historischen Befunde haben
engere Prädikate als der aktuelle Vertrag aus sieben Komponenten: Eine
lokalisierte Last ist kein quantifiziertes Budget, ein sichtbarer zweiter
Modus kein vollständiges E8-Modusaudit. Die korrigierte Zuordnung weist
F-R2a C4 (mit Abhängigkeit von C6), F-R2b und F-R6a C6 sowie F-R6b C7 zu;
F-R8 gehört zur getriggerten Erweiterung E8 außerhalb dieser Zuordnung.
Eine erklärte Normalisierungskonvention wird als teilweise liefernd (I)
eingestuft; sie ist keine Herleitung der Verträglichkeit zwischen
Containern. Wir korrigieren unsere bisherige allgemeine Behauptung, das
Volumen einer Menge dominiere den Betrag jedes Elements: Das zitierte
ATS-III-Argument betrifft eine Tensorprodukt-Box unter bestimmten
Integralitäts- und Gewichtshypothesen. Ein einfaches Beispiel einer
verschobenen Kugel widerlegt allein unsere allgemeine Paraphrase.
ATS-III-Boxeinschluss und die Normalisierungs- und Maßbrücke von ATS IV
bleiben getrennte Verpflichtungen. Spätere Klassifikationen der
Normalisierung und effektiven Gewichte werden innerhalb ihrer genannten
Hypothesen zitiert; sie schließen die indizierte Log-Volumen-, Hüllen-,
Maß- und Modusverträglichkeit nicht. Aus diesem Audit folgt keine
Behauptung über die Wahrheit von abc oder die Korrektheit von IUT und ATS.
\end{abstract}

\clearpage
\tableofcontents
\clearpage

%=====================================================================
\section{Auditauftrag, Quellenversionen und vorab festgelegte Reichweite}\label{sec:mandate}

Dieser Beitrag ist Teil~4 der abc-Gap-Serie. Teil~1~\cite{part1}
analysierte den Vergleich im Transportschritt von Korollar~3.12
aus IUTchIII~\cite{mochizukiIUT}. Teil~2~\cite{part2} definiert den
\emph{Bridge-Contract} (Brückenvertrag); Teil~3~\cite{part3} liefert
dessen zeilengenaues Ledger-Instrument. Beide sind veröffentlichte
Arbeitspapiere. Der vorliegende Beitrag wendet ausgewählte historische
Prädikate auf die exakten ATS-Versionen an
\cite{joshiI,joshiII,joshiIIhalf,joshiIII,joshiIV}.
Die Korrektur vom 2.~Oktober 2026 verarbeitet die inventarisierten
Zuordnungs-, Quellenparaphrase- und Statusimportfehler unseres
v0.2-Textes; sie führt kein neues Beweisaudit der gesamten Serie durch.

\subsection{Grundregeln}

Erstens \textbf{kein Autoritätsargument}: Weder \glqq Joshi behauptet, sie gelte\grqq{}
noch \glqq Scholze--Stix~\cite{scholzestix} behaupten, sie gelte nicht\grqq{} zählt als
Beleg; jedes Verdikt ist im wörtlichen Wortlaut der auditierten Quellen
verankert. Joshis eigene Statustabellen zu den Scholze--Stix-Einwänden werden
als Selbstauskunft behandelt, nicht als Beleg. Zweitens \textbf{kein
Wahrheitsverdikt}: Das Audit beurteilt, ob der Text vorab festgelegte
Vertragsklauseln bedient, nicht, ob abc wahr ist, ob Joshis Beweis korrekt ist
oder ob die Kritiker recht haben. Drittens \textbf{Statusdisziplin}: Die
auditierten Texte sind öffentlich zugängliche, versionierte arXiv-Preprints,
keine begutachteten Veröffentlichungen; ATS IV trägt die Kopfzeile
\glqq Preliminary version for comments\grqq{} und beginnt mit \glqq[This is still a work in
progress!]\grqq{}. Wir schreiben daher durchgängig \glqq öffentlich zugängliche
Fassung\grqq{}, binden die exakten Versionen (Tabelle~\ref{tab:manifest}) und
behandeln den Preprint-Status als Vorbehalt zur Versionsstabilität, der je
Befund vermerkt wird.

\subsection{Vorab festgelegte Reichweite: ein Vier-Zeilen-Test}\label{sec:scope}

Der Auditauftrag --- im Ledger-Instrument dieser Serie fixiert, bevor dieser
Test lief --- beschränkt den Test auf vier Ledger-Zeilen. \textbf{Dies ist ein
beauftragter Vier-Zeilen-Test, kein vollständiges Audit der Kernzeilen
F-R0--F-R7 mit getriggertem E8.} Die Zeilen R0,
R1, R3, R5, R7 waren nicht Teil des Auftrags (R5 wird nur insoweit berührt, als
R4 es erzwingt). Tabelle~\ref{tab:audit-scope} fixiert die einbezogenen und
ausgeschlossenen Zeilen.

\begin{table}[htbp]
\centering\small
\setlength{\tabcolsep}{4pt}
\begin{tabularx}{\textwidth}{@{}l >{\raggedright\arraybackslash}X >{\raggedright\arraybackslash}p{5.5cm}@{}}
\toprule
Zeile & Grund der Vorauswahl & Nicht abgedeckt \\
\midrule
R2 & numerischer Pilot-Transport; vergleichender Unterscheidungsfall & kein R0/R1-Audit \\
R4 & sichtbare Behandlung der $j^2$-Last & R5 nur über R4 \\
R6 & Vertragskomponenten 6--7; die zentrale Brückenfrage & kein Audit der Beweiskorrektheit \\
F-R8 & Sichtbarkeit des zweiten Modus & keine Vollständigkeit aller Hilfsmodi \\
\bottomrule
\end{tabularx}
\caption{Vorab festgelegte Reichweite des Vier-Zeilen-Audits. Die Tabelle
dokumentiert die Auftragsgrenze und verhindert, dass die Zeilenbefunde als
vollständiges Zertifikatsaudit gelesen werden.}
\label{tab:audit-scope}
\end{table}

\subsection{Operative Verdiktregeln}\label{sec:rules}

Jede atomare Behauptung erhält eines von drei Verdikten. \emph{Liefert} $=$ Das
vorab formulierte Minimalprädikat der Zeile wird vom zitierten Text bedient.
\emph{Liefert teilweise} $=$ Mindestens ein ausdrückliches Teilprädikat bleibt
unbedient. \emph{Liefert nicht} $=$ Das Minimalprädikat bleibt unbedient.
\textbf{Ein Bestehen auf Zeilenebene ist kein Bestehen des Zertifikats}: Es
bescheinigt allein das Minimalprädikat der Zeile, nie den vollständigen Vertrag
aus sieben Komponenten. Die Konfidenz wird in drei getrennten Dimensionen
berichtet: \textbf{Q} (Quellenbindung: Ist der wörtliche Anker eindeutig?),
\textbf{K} (Kriteriumspassung: Bedient der Wortlaut das vorab formulierte
Prädikat?), \textbf{S} (Substanzstatus: geprüft / nicht geprüft / extern als
offen markiert).

Operatives Glossar (die unten verwendeten Minimalpaare):
\emph{transportiert} $=$ eine numerische Größe zusammen mit einem
ausdrücklichen Transformationsgesetz; \emph{invariant} $=$ derselbe Ausgabewert
nach getypter Identifikation; \emph{adressiert} $=$ Last und ihre
kompensierende Stelle sind benannt; \emph{quantitativ verbucht} $=$
Betrag/Regel in einem vollständigen Budget bilanziert; \emph{sichtbarer Modus}
$=$ Modus, Rolle und betroffene Schranke ausdrücklich benannt.

\subsection{Zuordnung zum aktuellen Vertrag}\label{sec:crosswalk}

Der historische Auftrag bleibt R2/R4/R6 zuzüglich R8. Im aktuellen
Namensraum von Teil~2~\cite{part2} heißen die folgenden atomaren
Befunde F-R2a, F-R2b, F-R4a, F-R4b, F-R6a, F-R6b und F-R8.
F-R2a gehört zu C4 mit Abhängigkeit von C6; F-R2b und F-R6a zu C6;
F-R4a und F-R4b zu C5; F-R6b zu C7. Die getypte Auswertungsidentität
$\nu_C\circ i_q=\nu_q$ ist eine C6/F-R6a-Verpflichtung, keine
Komponente~1; ein sichtbares Skalierungsgesetz allein bescheinigt sie
nicht. F-R8 ist die getriggerte Erweiterung E8 außerhalb der
Zuordnung zu den sieben Komponenten. Ohne relevanten Mehrmodus-,
Vorzeichen- oder Ordnungstrigger ist sie X (von einem Nenner
ausgeschlossen); hier löst der ausdrückliche Frobenius-Shift sie aus.

Die historischen Adress- und Sichtbarkeitsprädikate bleiben als
solche erhalten. Ein historischer Befund \emph{liefert} für F-R4a oder
F-R6b behauptet weder das aktuelle C5-Budget noch die aktuelle
hypothesenweise C7-Abdeckung. Ebenso begründet F-R8-Sichtbarkeit
keine vollständige E8-Modusabdeckung. Wir geben daher weder eine
aktuelle Gesamtzertifikatswertung noch einen Anteil bestandener
Zeilen an. Nach der aktuellen Grammatik ist eine dokumentierte
Konvention höchstens I (liefert teilweise), kein E+ (liefert eine
Verträglichkeitsgleichung). Joshi ist ein vergleichender
Unterscheidungsfall; er ist nicht die akzeptierte Kontrolle der
gesonderten historischen v1-T-Illustration aus Teil~2.

\subsection{Quellenmanifest}\label{sec:manifest}

\begin{table}[htbp]
\centering\small
\setlength{\tabcolsep}{4pt}
\begin{tabularx}{\textwidth}{@{}l >{\raggedright\arraybackslash}X >{\raggedright\arraybackslash}p{5.5cm}@{}}
\toprule
ID & Quelle (exakt auditierte Version) & Rolle \\
\midrule
ATS III & arXiv:2401.13508v4 (revidiert 24.~Februar 2025)~\cite{joshiIII} & Konstruktion zu Korollar~3.12 (R2, R4, R6) \\
ATS IV & arXiv:2403.10430v2 (revidiert 24.~Februar 2025)~\cite{joshiIV} & abc-Schluss, Brückenstelle (R2, R6, F-R8) \\
ATS II($\tfrac12$) & arXiv:2305.10398v12~\cite{joshiIIhalf} & Normalisierungsdaten (Statushinweis A1) \\
\bottomrule
\end{tabularx}
\caption{Auditierte Quellen. Die Zeilenanker \glqq extract ll.\grqq{} beziehen sich auf
die Klartext-Auszüge genau dieser PDF-Versionen in den historischen
Auditunterlagen. Sie sind lokale Findehilfen, keine unabhängig
veröffentlichte Evidenz; die exakten Versions-, Abschnitts- und
Theoremangaben sind die öffentlichen Primärquellenanker. Auditdatum:
12. August 2026; Statushinweise vom 13./15. August 2026
(Abschnitt~\ref{sec:notices}).}
\label{tab:manifest}
\end{table}

\subsection{\texorpdfstring{Vier \glqq tautology\grqq{}-Stellen, auseinandergehalten}{Vier \textquotedbl tautology\textquotedbl -Stellen, auseinandergehalten}}\label{sec:tautology}

Die auditierten Texte verwenden das Wort \glqq tautology\grqq{} an vier Stellen mit
unterschiedlicher Reichweite; wer sie vermengt, gewinnt den falschen Eindruck,
Joshi nenne seinen eigenen arithmetischen Beweis eine Tautologie --- was er
nicht tut. (T-a) eine maßtheoretische Trivialität, die als Werkzeug dient
(ATS~III, Lemma~9.10.2.2(4)); (T-b) eine Aussage, die allein das
\emph{geometrische} Analogon betrifft (ATS~III, \S\,12.18); (T-c) eine Aussage
über Joshis eigene Variante $\tilde\Theta_{\mathrm{Joshi}}$, während das
fragliche Theorem für $\tilde\Theta_{\mathrm{Mochizuki}}$ bewiesen wird
(ATS~IV, \S\,1.6(4e)); und (T-d) die tragende Brückengleichheit im Beweis von
ATS~IV, Theorem~6.10.1.
\textbf{Allein T-d ist für dieses Audit tragend.}

%=====================================================================
\section{Das Ledger-Audit (atomare Befunde)}\label{sec:ledger}

\subsection{F-R2a --- Transport innerhalb des Links: liefert}\label{sec:r2a}

\emph{Minimalprädikat: Das numerische $q$-Readout wird entlang des Theta-Links
durch ein ausdrückliches numerisches Gesetz transportiert, wobei der
$j^2$-Faktor sichtbar mitgeführt wird.} ATS~III, Theorem~4.2.2.1(4), gibt die
Skalierungsrelation der Bewertungen
$|-|_{K'_{y_{w,j}}} = |-|^{j^2}_{K'_{y_{w,1}}}$ an den ungeraden semistabilen
Stellen an (\glqq Then one has the following relationship between their
valuations\grqq{}, Gl.~(4.2.2.2); extract ll.~1544--1560); \S\,9.9 führt den
$q$-Wert ausdrücklich hindurch (Gln.~(9.9.3)--(9.9.4); extract
ll.~6458--6484); und \S\,9.10.2 liefert die konkrete Log-Volumen-Skala (normiert
durch $\mathrm{LogVol}(O_{L'_w}) = 0$, $\mathrm{Vol}(\lambda O) = |\lambda|$;
extract ll.~6562--6596). Das ist \emph{Transport durch ein ausdrückliches
numerisches Gesetz}, nicht Invarianz eines Wertes: Bescheinigt wird eine
kovariante Transformationsregel mit konkreter numerischer Zielskala, die das historische Transportgesetz-Prädikat F-R2a/C4 (mit Abhängigkeit von C6) bedient. Die volle getypte Identität
$\nu_C \circ i_q = \nu_q$ ist in der Quelle nicht ausgeschrieben und wird hier
nicht bescheinigt.
\textbf{Verdikt: erfüllt das Zeilenprädikat. Q hoch / K hoch / S nicht
geprüft.}

\subsection{F-R2b --- Transport zwischen Containern: durch Konvention gesetzt}\label{sec:r2b}

\emph{Minimalprädikat: Die Gleichheit der Readouts über den globalen
Frobenius-Shift $y_0 \leftrightarrow \varphi(y_0)$ hinweg ist hergeleitet,
nicht gesetzt.} Hier wird die Gleichheit durch die Wahl \glqq the same global
normalization\grqq{} \emph{gesetzt} (ATS~IV, Thm.~6.10.1, Präambel; extract
ll.~3186--3189), mit einem Ein-Satz-Beweis (\glqq The equality \dots is a
consequence of the choice of the same global normalization \dots\grqq{}; extract
ll.~3262--3266), während sich die lokalen Bestandteile unter $\varphi$ ändern
und die Quelle selbst festhält, dass dies \glqq precludes direct comparison of
various local quantities\grqq{} (ATS~IV, Rem.~6.10.2; extract ll.~3208--3211). Kein
Widerspruch --- aber eine Konvention, keine hergeleitete Verträglichkeit.
\textbf{Verdikt: liefert teilweise (I): Eine Konvention ist dokumentiert;
hergeleitete Verträglichkeit bleibt unbedient. Q hoch / K hoch / S siehe
Abschnitt~\ref{sec:notices}.} Dies ist die aktuelle F-R2b/C6-Einstufung,
kein E+-Verträglichkeitsbefund.

\subsection{F-R4a --- die \texorpdfstring{$j^2$}{j\texttwosuperior}-Last: adressiert und zugewiesen}\label{sec:r4a}

\emph{Minimalprädikat: Die Last ist adressiert --- Ursprung, Ausbreitung und
kompensierende Stelle sind benannt; sie verschwindet nicht unverbucht in einer
Koordinate und kehrt nicht stillschweigend zurück.} Die Last entsteht allein an
den ungeraden semistabilen Stellen (Thm.~4.2.2.1, Gln.~(4.2.2.2)
gegenüber~(4.2.2.3); extract ll.~1544--1570); sie überlebt bis in die
numerische Auswertung (das $\sum_j j^2/\ell^{*2}$ aus \S9.9; extract
ll.~6476--6498); und ihre Rückkehr wird einer erzwungenen Renormierung an den
komplementären Stellen zugewiesen, mit der Produktformel als
Buchungsinstrument (Cor.~4.2.2.4, extract ll.~1600--1614; Thm.~4.6.1(4),
extract ll.~1786--1791). Das befürchtete Szenario --- stillschweigendes
Verschwinden und unverbuchte Rückkehr --- tritt im Text nicht ein.
\textbf{Verdikt: erfüllt das Zeilenprädikat. Q hoch / K hoch / S nicht
geprüft.}
Das historische Adressprädikat quantifiziert nicht alle Lasten,
Vorzeichen, Reste und kompensierenden Budgeteinträge des aktuellen C5;
daraus folgt kein aktueller C5/E+-Befund. Vermerkte Zitationsinkonsistenz: Die Quelle bezeichnet dieselbe Relation
abwechselnd als Theorem~4.2.2.1(3) und~(4); das Bezugsobjekt ist eindeutig.

\subsection{F-R4b --- quantitative Bilanz: offen}\label{sec:r4b}

\emph{Minimalprädikat: Die Kompensation ist quantitativ verbucht --- Betrag,
Eindeutigkeit und Budgetaufteilung sind gezeigt.} Die Existenz der
kompensierenden Renormierung wird in einem einzeiligen Beweis aus der
Produktformel-Nebenbedingung erschlossen (\glqq This is clear from
Theorem~4.2.2.1 and the global constraint placed by the requirement that the
product formula holds\grqq{}; extract ll.~1613--1614).
Die Restlast ist \emph{lokalisiert, aber nicht quantifiziert};
Vertragskomponente 5 (separiertes Unbestimmtheitsbudget) wird qualitativ
bedient, nicht quantitativ.
\textbf{Verdikt: liefert teilweise. Q hoch / K hoch / S nicht geprüft.}

\subsection{F-R6b --- Ordnungsrelation, Formebene: ausgeschriebener Beweis mit
adressierbarem Kernschritt}\label{sec:r6f}

\emph{Minimalprädikat: Die Ordnungsrelation ist als Theorem formuliert, mit
einem ausgeschriebenen Beweis, dessen Schritte adressierbar sind.} ATS~III,
Theorem~9.11.1 (extract ll.~6794--6806; in Mochizukis Notation Cor.~9.11.1.1,
extract ll.~6882--6888), formuliert die untere Volumenschranke; die Auswertung
existiert im Container als adelisches Log-Volumen, berechnet als Summe der
lokalen Beiträge (\glqq one can compute log Vol (\dots) by taking the sum \dots of
the local contributions at all primes\grqq{}; ATS~IV, \S\,1.5, extract ll.~764--769);
dies ist die in Abschnitt~\ref{sec:r6cc} geprüfte Auswertung von
Vertragskomponente 6, nicht Teil des hier festgehaltenen Befunds zur
Ordnungsrelation. Der zitierte Mechanismus ist quellenspezifisch:
ATS~III, Lemma~9.10.7.1 (PDF S.~125), behandelt eine Obermenge $S'$
einer Tensorprodukt-Box mit integralen $\tau_i\in O_{E_i}$,
Bestandteilen $s_i\in\tau_i O_{E_i}$ und Gewichten
$0<\gamma_i\le1$. Theorem~9.11.1 (PDF S.~127) verwendet den
behaupteten Einschluss dieser Box in die Konstruktion.
Keiner der beiden Sätze ist die allgemeine Aussage, das
Volumen einer beliebigen Menge beschränke den Betrag jedes Elements.
Unsere entsprechende v0.2-Paraphrase war falsch: Für das normierte
additive Haarmaß $\mu(\mathbb Z_p)=1$ liefert die verschobene Kugel
\[
 U=1+p\mathbb Z_p,\qquad 1\in U,\qquad
 \mu(U)=p^{-1}<|1|_p=1
\]
ein Gegenbeispiel. Es widerlegt unsere uneingeschränkte Paraphrase,
nicht Joshis Tensor-Box-Lemma, Theorem~9.11.1, IUT oder abc.
Insbesondere darf die Hypothese eines zentrierten Moduls aus ATS~III,
Lemma~9.10.2.2 (PDF S.~123), nicht stillschweigend entfallen.

Der historische Formbefund beschränkt sich auf ein formuliertes
Theorem und einen adressierbaren Beweis. Der Nachweis, dass die
tatsächlich konstruierte Theta-Menge die Box, Integralstruktur und
anwendbaren Gewichtshypothesen liefert, bleibt die getrennte
A2-Verpflichtung. Benannte Beweisadressen allein erfüllen das aktuelle
C7 nicht: Jede tragende Hypothese, einschließlich der
Einschlusshypothese, benötigt einen eigenen anwendbaren Anker.
Die Verfügbarkeit adelischer Auswertung allein begründet keine
C6-Verträglichkeit.
\textbf{Verdikt: erfüllt das Formprädikat. Q hoch / K hoch / S extern als
offen markiert (A2).}

\subsection{F-R6a --- Verträglichkeit zwischen Containern: gesetzt und nach der
Selbstbeschränkung der Quelle nur global}\label{sec:r6cc}

\emph{Minimalprädikat (Vertragskomponente 6, globale Form): Die verbindende
Verträglichkeit zwischen den beiden Auswertungsrouten ist hergeleitet;
Vertragskomponente 7 nutzt sie ihrerseits, um unter zulässiger Verknüpfung
$\nu_q(\mathrm{Obj}_q) \in H_\Theta$ zu folgern.} Für den abc-Schluss wird die
untere Schranke mit $\varphi(y_0)$ berechnet, die obere mit $y_0$; verbunden
werden sie durch die mittlere Gleichheit von ATS~IV, Theorem~6.10.1, deren
vollständiger Beweis lautet: \glqq The middle equality is a tautology using the
fact that the number fields provided by $y_0'$ and $y_0$ are identically
normalized\grqq{} (T-d; extract ll.~3272--3273). Zwei Klarstellungen halten den
Standard ehrlich. Erstens ist die Ordnungsklausel des Zertifikats
\emph{global}; die Selbstbeschränkung der Quelle --- man \glqq cannot claim that
$z_p \le z_p'$ holds for all primes $p$ without violating the middle
equality\grqq{} (ATS~IV, \S\,1.5, extract ll.~816--819) und \glqq one cannot make local
term by term comparisons at each prime\grqq{} (extract ll.~772--775) --- ist eine
\emph{Reichweitengrenze der Behauptung}, kein verletztes Kriterium. Zweitens
bleibt auf der Formebene die \emph{Herleitung} der Verträglichkeit zwischen den
Containern unbedient: Sie wird durch definitorische Normalisierung gesetzt, und
zu den verbleibenden Verpflichtungen gehören der indizierte
E2-Log-Volumenvergleich, Umgebungs- und Integralstrukturen,
Theta-Mengen- und Hüllenverträglichkeit sowie Maß- und
Gewichtskonventionen (Abschnitt~\ref{sec:notices}, Hinweis A1).
Die Normalisierungsidentifikation auf Bewertungsebene allein
trägt diese Verpflichtungen nicht.
\textbf{Verdikt: liefert teilweise (erklärte Konvention; Herleitung
unbedient). Q hoch / K hoch / S siehe Hinweise.}

\subsection{F-R8 --- der zweite Modus: sichtbar und erklärt}\label{sec:r8}

\emph{Minimalprädikat: Der von der Brücke verwendete zweite Modus ist sichtbar
und seine Rolle erklärt (nicht: ein vollständiges Inventar aller Hilfsmodi).}
Der zweite Modus ist der globale Frobenius-/Log-Shift, das Paar
$(y_0, y_0' = \varphi(y_0))$. Er wird benannt; seine Wirkung wird der Richtung
nach charakterisiert (die Beträge der Tate-Parameter steigen unter $\varphi$;
ATS~IV, \S\,1.4(2), extract ll.~567--569); sein asymmetrischer Gebrauch wird
erklärt (\glqq the lower bound will be computed using the arithmeticoid
$y_0'=\varphi(y_0)$ while the upper bound will be computed using $y_0$ (the
upper bound is not independent of this choice)\grqq{}; \S\,1.5, extract
ll.~733--743); er wird geometrisch gedeutet (Steigungsvergleich
Frobenius-verschobener Bündel; Rem.~6.10.3, extract ll.~3213--3235); und er
wird als Korrektur einer früheren eigenen Fassung des Beitrags kenntlich
gemacht (\glqq[The previous release of this paper had omitted this
Frobenius-shifting.]\grqq{}; extract l.~755). \textbf{Verdikt: erfüllt das
Sichtbarkeitsprädikat. Q hoch / K hoch / S nicht geprüft.}

Modus-Inventar (nur Sichtbarkeitsstatus; Vollständigkeit wird \emph{nicht}
behauptet): der Frobenius-/Log-Shift (sichtbar, verwendet); die abzählbar
unendliche Wahl des Basispunkts $y_0$ (sichtbar, unverbucht; ATS~III,
Rem.~4.4.2, extract ll.~1723--1731); die beabsichtigte, aber nicht verfügbare
Mittelung über Arithmetikoide (sichtbar, nicht verfügbar; ATS~IV,
\S\,4.5(4)/(5): \glqq one is far from being able to do this\grqq{}, extract
ll.~2100--2121); die nur über ihre Summen eingeschränkten Hilfsgrößen
$z_p, z_p'$ je Stelle (sichtbar, unverbucht; ATS~IV, Gln.~(1.5.4)/(1.5.5),
extract ll.~798--819).

%=====================================================================
\section{Externe Statushinweise: korrigierte Reichweite}\label{sec:notices}

Der historische Audit erfolgte am 12.~August 2026; die Hinweise
A1--A3 wurden am 13./15.~August vermerkt. Ihre früheren abschließenden
Deutungen werden hier zum 2.~Oktober 2026 korrigiert. Diese Hinweise
sind quellenberichtete Folgeeinschätzungen, keine durch diesen
Formaudit begründeten Beweise. Die Teile~5 und~6 sind inzwischen
veröffentlichte Arbeitspapiere \cite{part5,part6}; öffentliche
Verfügbarkeit ist keine mathematische Abnahme und keine
Aufwertung ihrer Behauptungen.

\begin{table}[htbp]
\centering\footnotesize
\begin{tabularx}{\textwidth}{@{}l >{\raggedright\arraybackslash}X >{\raggedright\arraybackslash}X@{}}
\toprule
Hinweis & Korrigierte Reichweite & Wirkung auf dieses Audit \\
\midrule
A1 (13. Aug.) & Normalisierung auf Bewertungsebene; indizierte Hüllen-/Maß-/Gewichts- und E2-Verträglichkeit offen & F-R2b und F-R6a bleiben teilweise \\
A2 (13. Aug.) & ATS-III-Tensor-Box-Einschluss und anwendbare Hypothesen offen & historischer Beweisadressbefund bleibt; keine C7-Substanzzertifizierung \\
A3 (15. Aug.) & bedingte Klassifikation effektiver Gewichte unter fixierten Körper-/Stellen- und Identitätshypothesen & kein unbedingtes F-R6a und kein globaler Zweigabschluss \\
\bottomrule
\end{tabularx}
\caption{Korrigierte Lesart der historischen Hinweise, 2.~Oktober 2026.
Die Daten kennzeichnen ihre Herkunft; der aktuelle Text begrenzt ihre Übernahme.}
\label{tab:status-notices}
\end{table}

A1 betrifft die mittlere Gleichheit in ATS~IV, Theorem~6.10.1
(PDF S.~66--67), mit den Reichweitengrenzen aus \S\,1.5 und
Remarks~6.10.2--3. Die korrigierte Normalisierungsbeilage
v1.2~\cite{gnc} unterscheidet eine Identifikation auf Bewertungsebene
von den weiteren Entscheidungen für den Vergleich indizierter
E2-Log-Volumina. Erklärte Normalisierung klärt weder Umgebungs- und
Integralstruktur noch Theta-Mengen-, Hüllen-, Maß-, Gewichts- oder
Ordnungsverträglichkeit. Ein vorgeschlagener Frobenius-Mechanismus
zur Mengenstabilität bleibt quellenberichtet; er ist kein
nachgewiesener E2-Transportsatz. F-NEU-1 bleibt offen.

A2 betrifft ATS~III, Lemma~9.10.7.1 und Theorem~9.11.1
(PDF S.~125 und~127). Seine Box-Einschlussverpflichtung ist von
A1s Verpflichtung zur Normalisierung und zum Maß in ATS~IV getrennt. Die Korrektur
unserer allgemeinen Volumenparaphrase erfüllt A2 nicht und widerlegt
das quellenspezifische Lemma nicht. Die aktuell zitierte
Einschätzung aus Teil~5~\cite{part5} bleibt als datierter Bericht
erhalten; sie wird in diesem Beitrag nicht unabhängig zertifiziert.

A3 ist auf Größen begrenzt, deren Faktorisierung durch effektive
Produktformelgewichte gezeigt ist. Die berichtete Klassifikation
erfordert einen fixierten Zahlkörper, dieselbe vollständige
Stellenmenge, reelle effektive Gewichte und eine Identität für alle
von null verschiedenen Elemente unter den genannten Betragskonventionen;
quellenspezifische punktweise Normalisierungsidentitäten stellen
weitere Bedingungen. Die volle Produktformel allein lässt einen
gemeinsamen Skalar übrig; eine punktweise Normalisierungsidentität
ist eine stärkere Hypothese. Keine der beiden Folgerungen gilt
automatisch für Log-Volumina, Maße, Theta-Mengen oder Hüllen.
Die Alle-Stellen-Hypothese darf eine rein nichtarchimedische
Stellenreichweite nicht stillschweigend ersetzen. Die
Teile~5 und~6~\cite{part5,part6} benötigen ihre eigene
quellengebundene Prüfung; weder ein unbedingtes Verträglichkeitsverdikt
zwischen Containern noch ein endgültiger Abschluss des Lesartzweigs
(c) oder ein positiver E2-Satz wird hier übernommen.

%=====================================================================
\section{Synthese (konservativ)}\label{sec:synthesis}

Tabelle~\ref{tab:synthesis} folgt auf die korrigierten Hinweise
und unterscheidet historische Adressbefunde von aktuellen
Komponentenverpflichtungen (Abschnitt~\ref{sec:crosswalk}).

\begin{table}[htbp]
\centering\footnotesize
\begin{tabularx}{\textwidth}{@{}l >{\raggedright\arraybackslash}X >{\raggedright\arraybackslash}X@{}}
\toprule
Befund & Sichtbare Form & Nicht bescheinigte Verpflichtung \\
\midrule
F-R2a / C4 & ausdrückliches $j^2$-Transportgesetz & C6-Abhängigkeit; getypte Auswertungsidentität \\
F-R2b / C6 & erklärte Konvention; I & hergeleitete Verträglichkeit zwischen Containern (A1) \\
F-R4a / C5 & Last und Stelle adressiert (historisches Prädikat) & vollständiges quantifiziertes Budget \\
F-R4b / C5 & qualitative Lokalisierung; teilweise & Betrags-, Aufteilungs-, Rest- und Vorzeichenbuchung \\
F-R6b / C7 & Theorem und Beweisadressen (historisches Prädikat) & anwendbare Hypothesenanker; tatsächlicher Boxeinschluss (A2) \\
F-R6a / C6 & Auswertung verfügbar; globale Reichweite benannt; teilweise & getypte Auswertungsidentität; indizierte E2-Hüllen-/Maß-/Gewichtsverträglichkeit (A1) \\
F-R8 / E8 & Frobenius-Modus sichtbar (historisches Prädikat) & vollständige Abdeckung/Ausschließung ordnungsrelevanter Modi \\
\bottomrule
\end{tabularx}
\caption{Konservative Synthese ausgewählter Zeilen; keine aktuelle
Gesamtzertifikats- oder Beweisgütewertung. Q/K-Quellen- und
Kriteriumskonfidenz unterscheidet sich von ungeprüfter oder
offener mathematischer Substanz.}
\label{tab:synthesis}
\end{table}

Die Quelle liefert identifizierbare Skalierungs-, Kompensationsstellen-
und Modusaussagen. F-R2b und F-R6a bleiben teilweise; F-R4b bleibt ein
lokalisiertes, aber nicht quantifiziertes Budget. Die historischen
F-R4a-, F-R6b- und F-R8-Befunde betreffen engere Prädikate als
das aktuelle C5, C7 und E8. Diese Beobachtungen erlauben die
Prüfung benannter Verpflichtungen, kein rückwirkendes Bestehen
des Gesamtvertrags. Das Gegenbeispiel der verschobenen Kugel
korrigiert allein die eigene Paraphrase dieses Beitrags. Die
Verpflichtungen A1 und A2 bleiben getrennt; die bedingte
A3-Klassifikation schließt keine der beiden.

\begin{table}[htbp]
\centering\footnotesize
\begin{tabularx}{\textwidth}{@{}l >{\raggedright\arraybackslash}X >{\raggedright\arraybackslash}X@{}}
\toprule
Ebene & Aktuelle Reichweite & Verbleibende Last \\
\midrule
Normalisierung & Identifikation/Konvention auf Bewertungsebene & getypter/indizierter Vergleich nicht nachgewiesen \\
Maß / E2 & keine allgemeine Positivmaß-Klassifikation übernommen & Umgebungs-, Hüllen-, Maß-, Gewichts- und Ordnungsverträglichkeit \\
Mengenstabilität & möglicher Frobenius-Mechanismus quellenberichtet & kein nachgewiesener indizierter E2-Mengen-/Hüllentransport \\
Box / A2 & konkrete Tensor-Box-Hypothesen und Einschluss & tatsächliche Konstruktion und anwendbare Hypothesenanker \\
Effektive Gewichte / A3 & bedingte Klassifikation für fixierten Körper/alle Stellen & Anwendbarkeit, Stellenreichweitenabgleich; keine Maßfolgerung \\
\bottomrule
\end{tabularx}
\caption{Aktuelle Belastungskarte, 2.~Oktober 2026. Die früheren
unbedingten und abschließenden Übernahmen werden zurückgenommen;
offene Verpflichtungen sind keine neuen Widerlegungen von ATS.}
\label{tab:burden}
\end{table}

%=====================================================================
\clearpage
\section{Grenzen der Behauptung}\label{sec:claims}

Dieses Audit erhebt keinen Anspruch über die Wahrheit der abc-Vermutung, die
Korrektheit von IUT, die Korrektheit der ATS-Serie oder die letztliche
Reparierbarkeit des Transportschritts. Es beurteilt öffentlich zugängliche
Preprint-Fassungen an einer vorab festgelegten Vertragsform, mit wörtlichen
Ankern; ein unbedientes Zeilenprädikat lokalisiert eine Last, es widerlegt keinen
Satz. Das Audit ist auf die vier beauftragten Zeilen beschränkt; weder
Vollständigkeit des Modus-Inventars noch ein Bestehen des vollständigen
Zertifikats noch eine Aggregation der Zeilenverdikte zu einer
Beweisgütewertung wird behauptet. Die Substanzanalyse und Reparaturklassifikationen sind Gegenstand
der veröffentlichten Arbeitspapiere Teile~5 und~6
\cite{part5,part6}; dieser Beitrag übernimmt allein ausdrücklich
begrenzte, datierte Quellenberichte, keine neue Abnahme ihrer Mathematik.

% Lokale, begrenzte Anpassung von ai_disclosure_STANDARD_de, 2026-10-02.
\section*{KI-Offenlegung}
\noindent
KI-Systeme wurden für Forschungsdiskussion, Literatur- und Quellenarbeit,
mathematische Analyse, Text, Übersetzung und Satz eingesetzt. Frühere
KI-gestützte Einschätzungen unterliegen weiterhin den in diesem
Arbeitspapier genannten Quellen- und Behauptungsgrenzen.
\medskip
\noindent
Die Qualitätssicherung dieser Korrekturfassung umfasst getrennte
quellengebundene mathematische und zweisprachige Änderungsprüfungen,
die Kontrolle von Quellenversionen und bibliografischen Identitäten,
Kompilierung und Sichtprüfung beider PDFs. Diese Kontrollen betreffen
die genannten Manuskriptkorrekturen; sie zertifizieren die
ATS-Beweiskette nicht. Weder ein neues rechnerisches Experiment
noch eine Beweisprüfung der gesamten Serie wird berichtet.


%=====================================================================
\begin{thebibliography}{99}
\small
\raggedright
\sloppy
\setlength{\itemsep}{2pt plus 0.5pt minus 0.5pt}

\bibitem{joshiI}
K.~Joshi,
\emph{Construction of Arithmetic Teichmüller Spaces I},
arXiv:2106.11452v4 [math.AG], 2021 (revidiert 24.~Februar 2025).

\bibitem{joshiII}
K.~Joshi,
\emph{Construction of Arithmetic Teichmüller spaces II: Proof of a local
prototype of Mochizuki's Corollary 3.12},
arXiv:2303.01662v3 [math.AG], 2023 (revidiert 24.~Februar 2025).

\bibitem{joshiIIhalf}
K.~Joshi,
\emph{Construction of Arithmetic Teichmüller Spaces II$\frac{1}{2}$:
Deformations of Number Fields},
arXiv:2305.10398v12 [math.AG], 2023 (revidiert 24.~Februar 2025).

\bibitem{joshiIII}
K.~Joshi,
\emph{Construction of Arithmetic Teichmüller Spaces III: A `Rosetta Stone'
and a proof of Mochizuki's Corollary 3.12},
arXiv:2401.13508v4 [math.AG], 2024 (revidiert 24.~Februar 2025).

\bibitem{joshiIV}
K.~Joshi,
\emph{Construction of Arithmetic Teichmüller Spaces IV: Proof of the
abc-conjecture},
arXiv:2403.10430v2 [math.AG], 2024 (revidiert 24.~Februar 2025).

\bibitem{mochizukiIUT}
S.~Mochizuki,
\emph{Inter-universal Teichmüller Theory I: Construction of Hodge
Theaters}, Publ.\ Res.\ Inst.\ Math.\ Sci.\ \textbf{57} (2021), no.~1/2,
3--207;
\emph{II: Hodge--Arakelov-Theoretic Evaluation}, ibid., 209--401;
\emph{III: Canonical Splittings of the Log-Theta-Lattice}, ibid., 403--626;
\emph{IV: Log-Volume Computations and Set-Theoretic Foundations},
ibid., 627--723.

\bibitem{scholzestix}
P.~Scholze und J.~Stix,
\emph{Why abc is still a conjecture},
Unveröffentlichtes Manuskript, Fassung vom 16.~Juli 2018.
Verfügbar unter
\url{https://www.math.uni-bonn.de/people/scholze/WhyABCisStillaConjecture.pdf}

\bibitem{part1}
L.~Geiger,
\emph{The Comparison in IUTchIII, Corollary 3.12: From Attempted Repair
to Structural Diagnosis},
Zenodo, 2026, version v6.4.
DOI: \nolinkurl{10.5281/zenodo.23081604}.

\bibitem{part2}
L.~Geiger,
\emph{The Bridge-Contract Criterion: A Form Audit for Transport Arguments,
with an Instantiation in the abc Literature},
Zenodo, 2026, version 0.2.
DOI: \nolinkurl{10.5281/zenodo.21960477}.

\bibitem{part3}
L.~Geiger,
\emph{The Bridge Certificate: A Line-Level Audit Instrument for
Inter-Theoretic Transfer Arguments, with Four Applications to the abc Literature},
Zenodo, 2026, version 0.3.
DOI: \nolinkurl{10.5281/zenodo.23090371}.

\bibitem{part5}
L.~Geiger,
\emph{The Joshi ATS Series under the Bridge-Certificate Standard:
A Targeted Substance Audit and an All-Place Rigidity Classification
of Effective Normalization Weights},
Zenodo, 2026, version 0.2.
DOI: \nolinkurl{10.5281/zenodo.21956399}.

\bibitem{part6}
L.~Geiger,
\emph{What the Product Formula Allows: A Rigidity Theorem and an
Effective-Weight Test for the ATS Height Subchain},
Zenodo, 2026, version 0.2.
DOI: \nolinkurl{10.5281/zenodo.21952487}.

\bibitem{gnc}
L.~Geiger,
\emph{Global Normalization in Joshi's Arithmetic Teichmüller Theory:
What the Identification Carries --- and What Carries the Theorem (DRAFT)},
Zenodo, 2026, version 1.2.
DOI: \nolinkurl{10.5281/zenodo.23092166}.

\bibitem{landscapeA}
L.~Geiger,
\emph{The abc Landscape: Obstructions, Failed Routes, and HCT Diagnostics
(DRAFT)},
Zenodo, 2026.
Concept DOI: \nolinkurl{10.5281/zenodo.21916900}
(zitierte Version 1.1, 13.~August 2026, DOI: \nolinkurl{10.5281/zenodo.21924338}).

\end{thebibliography}

\end{document}
